% Edition 22: French physical 88, paragraph 2, through 90 before section15.
% English physical 57 More delicate... through 59 before section15.
锥投影还有更精细的性质，即它的微分性质，尤其是在$X$在$k$上光滑、$f:X\to P$非分歧的情形；第12节研究的正是这些性质。回顾一下：在$Q(C)$的某个余维为1的子集$Z$之外，定义在$\widetilde X(C_\eta)$上的态射$p_{C_\eta}$是光滑的。而第12节所勾勒的更细致分析表明（若我们尚未作出这一分析，则将会表明）：如果$X$的各不可约分支的维数均$\geq m$，那么在$Q(C_\eta)$的某个余维$\geq2$的子集$Z'\subset Z$之外，纤维
\[
p_{C_\eta}^{-1}(\xi)=X^{(m)}_\xi
\]
至多只有几何意义下的普通奇点；而且，若$f$是浸入且$X$几何不可约，则这样的点至多一个，该点于是必为$k(\xi)$上的有理点。这些断言至少在$k$的特征为0时成立；或者，如第9节那样，以$\varphi_n f$（$n\geq2$）代替$f$，在这个条件下也成立。\cmt{V22-DIF001}{法88明确保留“表明（或者将会表明，若我们尚未作）”的草稿语气；此段仍是作者拟议分析的回顾，不因本次译校而被升级为已完整补证的定理。普通奇点、唯一性及有理性所依赖的第12节结论仍有当前未决表所列界限。法稿第一次写$Q(C)$，随后写$Q(C_\eta)$；前者按本节泛中心约定理解。纤维式中原省去的中心下标在此明确为$C_\eta$。}

也应当给出$\dim X\leq m$、从而$p_{C_\eta}$定义在$X_{k(\eta)}$上时的微分性质。我这里只指出下列性质；证明应该容易，留给Dieudonné。\cmt{V22-DIF002}{法88手改可读为“doit être facile et est laissée à Dieudonné”；英文另加“或Blass”并标$[\mathrm{Tr}]$，这是英文编者的插语，不能当作法稿原句。下列命题在原稿中没有展开证明；本版的独立核查另列于报告。}

\egas{14.4}{命题}
设$f:X\to P$非分歧，且$\dim X\leq m$。设$T$是$X$的一个在$k$上有限的子预概形。则：

\noindent(a) 若$f$是浸入且$m\geq1$，则$p_{C_\eta}$在$T_{k(\eta)}$上的限制是根态射，即“几何单射”。\cmt{V22-DIF003}{法88原打字“un monomorphisme”已划去，手改为“radiciel i.e. géométriquement injectif”，故这里必须用根态射（普遍单射），不能恢复成已经被作者删除的单态射。原稿没有排除本节允许的$m=0$；若$X=T$是$\mathbb P^1$中两个不同的$k$有理点，泛投影到$\mathbb P^0$就合并两点，原断言为假。正文新增$m\geq1$这一充分限定；报告另说明$m=0$时应要求$T$至多有一个几何点。有限子预概形可以非约化，普遍单射不等于单态射。}
若再设$Y$为$X$的一个维数$\leq m-1$的闭子集，并且$Y\cap|T|=\varnothing$，则\cmt{V22-DIF004}{法稿和英文都漏写$Y\cap|T|=\varnothing$。取$Y=T$为直线上的一个有理点，即可否定原来的无条件空交式。正文补入这个必要条件；报告V22-DIF-E02用割线与中心的关联维数计数证明，在$f$为浸入和原维数限制下，它也是充分的。此处是集合的空交；若把$T$保留为子预概形，所对应逆像子预概形同样为空。}
\[
p_{C_\eta}^{-1}\bigl(p_{C_\eta}(Y_{k(\eta)})\bigr)
\cap |T_{k(\eta)}|=\varnothing.
\]


\noindent(b) 若$X$在$T$的各点光滑，则$p_{C_\eta}$在$T_{k(\eta)}$的所有点非分歧。\cmt{V22-DIF005}{法88末行手插可恢复为“et même en les points de $p_{C_\eta}^{-1}p_{C_\eta}(T_{k(\eta)})$”，即“甚至在该整条逆像的各点”。英文将连接语标为“字迹不清，询问A.G.”；现在文字可读，但这个加强断言在所列假设下是错的。报告V22-DIF-E03以一条光滑直线和另一条异面直线的非约化加厚的并给出反例：$T$处光滑并不能控制另一分支上、位于相同投影纤维中的非分歧性。故正文保留已证的$T_{k(\eta)}$上结论，把原加强句与公式完整记在本注；不能把“已辨读”与“已证明”混同。}

\egas{14.5}{命题}
设$\dim X\leq m-1$，$f:X\to P$为浸入，并且$X$在$k$上可分。设$Y_\eta$为$X_{k(\eta)}$在$Q(C_\eta)$中的概形论像。则诱导态射
\[
p'_{C_\eta}:X_{k(\eta)}\longrightarrow Y_\eta
\]
是双有理的。\cmt{V22-DIF006}{法89确实写“séparable”，没有写“光滑”；保留原术语。报告对双有理部分明确使用有限型、几何约化的解释，并逐个处理几何不可约分支。目标是概形论像，法稿用带撇号的$p'_{C_\eta}$；英文省去撇号，容易与向整个$Q(C_\eta)$的态射混淆。}
而且，对$X_{k(\eta)}$中位于$X$的某个光滑闭点之上的每一点$x$，$p'_{C_\eta}$在$x$处以及在
\[
(p'_{C_\eta})^{-1}\bigl(p'_{C_\eta}(x)\bigr)
\]
的各点处都是平展的。\cmt{V22-DIF007}{法89手插及英文均把此句写成“位于$X$的某个闭点之上”，没有“光滑”限定；这不是英文误读。该无条件断言不成立：$\mathbb P^3$中仿射单项式曲线$\Spec k[t^3,t^4,t^5]$的射影闭包，在原点的嵌入维数是3；当$m=2$时，其投影像是$\mathbb P^2$的子概形，切空间维数至多2，因此诱导态射在该点不可能非分歧，更不可能平展。正文加入“光滑”作为一个已补证的充分条件；报告V22-DIF-E04、E05分别给出反例和有限投影的局部同构证明。没有把这个额外条件说成原稿本来已有。}

注意以下推论：

\egas{14.6}{推论}
设$X$为无限域$k$上的一个不可约且可分的$n$维射影代数概形。则存在一个从$X$满射到$\mathbb P^{n+1}_k$中某个超曲面的双有理态射。\cmt{V22-DIF008}{法89删除“是双有理同构于”，手改为“存在从$X$到……上的双有理态射”；英文在此已正确反映修改。本版保持态射与有理映射的区别。报告V22-DIF-E06补出选取$k$上中心的非空开条件、有限性和像为超曲面的理由；它不依赖14.5中被反例否定的奇异闭点平展性。}

即使$X$是$P$的一个光滑、几何不可约、维数为$m-1=n$的闭子预概形，也不要认为锥投影$p_{C_\eta}$必定是浸入。事实上，若$k$无限，这就意味着存在一个在$k$上有理的中心$C$，具有同样的性质；于是$X$同构于$\mathbb P^{n+1}$中的一个非奇异超曲面。但即使$n=1$，即$X$是代数闭域上的光滑连通射影代数曲线，也容易构造不能嵌入$\mathbb P^2$的$X$的例子。\cmt{V22-DIF009}{法89的$p_{C_\eta}$及“必定是一个浸入”可辨，英文在此附“字迹不清”不再保留为法稿缺文。报告给出亏格2曲线与光滑平面曲线亏格公式的具体例证，明确这是译者补例。不能把双有理满射误加强为同构或闭浸入。}

同样，对于14.4，不要把已给出的陈述与下列一般不成立的断言混为一谈：$p_C$本身是单态射（前面的反例；或者，在$X$光滑且维数为$m$时，还有更明显的反例），或者$p_C$应当非分歧。\cmt{V22-DIF010}{法89清楚写“ou contre-exemple encore plus évident si $X$ est lisse de dimension $m$”。英文丢失“或者还有更明显的反例”，使$m$看似与前文$m-1$冲突。现在恢复完整连接，不再把这一处列为作者维数矛盾。}
对于后一项，可取$X$为一个光滑、不可约、维数为$m$的闭子概形，譬如假设$k$代数闭。若存在一个非分歧态射
\[
X\longrightarrow Q\simeq\mathbb P^m,
\]
则由维数的理由，它将是平展的；但可以证明（见第VIII章！），这意味着
\[
X\xrightarrow{\sim}\mathbb P^m
\]
（因为$\mathbb P^m$是“单连通”的）。\cmt{V22-DIF011}{法稿括注为“sur $k$ alg clos disons”，可译“譬如假设$k$代数闭”，没有不可辨的数学假设。此处$X$承接为射影空间的闭子概形：非分歧给出准有限，固有性使它有限；光滑等维且目标正则给出平坦，从而是有限平展态射。还须用连通性及射影空间的代数基本群结论才能得到同构，不能把结论推广到任意非固有平展态射。原“第VIII章！”保留原引用身份，报告另用曲线例子独立验证一般非分歧断言确实不成立。}

14.4的直观几何意义是：$p_{C_\eta}$的分歧集在$k$上是“可变的”；更确切地说，当$\xi$在$\operatorname{Grass}^{m+1}(\bar k)$的一个开集中变化时，$p_{C_\xi}$的分歧集在$X(\bar k)$中变化，并且没有任何“固定点”……\cmt{V22-DIF012}{法89手插“在$k$上”及“固定点”可读。这里承接本段$X$光滑的情形；若只要求14.4(b)中$T$处光滑，不能据此宣称整个$X$的分歧轨迹没有固定部分，V22-DIF-E03的另一分支正提供反例。参数空间省写$P$沿用原稿的约定。}

当然，为了在本节中说明从泛点$\eta$到$\operatorname{Grass}^{m+1}(P)$的邻近点的过渡，并在需要时重新采用7.1的一般说明，必须考察下列图式：\cmt{V22-DIF013}{法90图中左下角清楚为$S$；英文物理58页误作$X$。旧稿虽已据基变换订正，但本轮核对法稿后，可以明确这是英文转写错误，不是原作者画错了图。}
\[
\begin{tikzcd}
X \arrow[d] & \widetilde X(C) \arrow[l] \arrow[d] \\
S & Q(C) \arrow[l]
\end{tikzcd}
\]

这些图式由各个$C\in\operatorname{Grass}^{m+1}(P)(S)$得到；更一般地，还要考虑经基变换$T\to S$之后，由$\operatorname{Grass}^{m+1}(P)$的$T$值点$\xi$得到的图式：
\[
\begin{tikzcd}
X_T \arrow[d] & \widetilde X(C_\xi)=\widetilde X_T(C_\xi) \arrow[l] \arrow[d] \\
T & Q(C_\xi) \arrow[l]
\end{tikzcd}
\]
它们是由下面的普遍图式经基变换
\[
\xi:T\longrightarrow \mathcal G,
\qquad\mathcal G=\operatorname{Grass}^{m+1}(P)
\]
得到的。这个普遍图式对应于$\operatorname{Grass}^{m+1}(P)$的典范$\mathcal G$值点：\cmt{V22-DIF014}{法90用带下划线的$T$与$C$区别普遍对象和试验对象；英文丢失下划线，又混入无所指的$J$，故看似把任意$T$等同于Grassmann概形。本版将普遍底改记$\mathcal G$、普遍中心改记$\mathcal C$，保留其区别。第二图右上等号后的波浪号在法稿中明确存在，英文漏掉；此处已恢复。}
\[
\begin{tikzcd}
X_{\mathcal G} \arrow[d] & \widetilde X(\mathcal C) \arrow[l] \arrow[d] \\
\mathcal G & Q(\mathcal C) \arrow[l]
\end{tikzcd}
\]
这里$\mathcal C$是$P_{\mathcal G}$的典范线性子簇。于是，上述态射
\[
\widetilde X(C_\eta)\longrightarrow Q(C_\eta)
\]
正是后一图式中$\mathcal G$态射
\[
\widetilde X(\mathcal C)\longrightarrow Q(\mathcal C)
\]
在泛纤维上诱导的态射；而泛纤维态射所具有的每一个可构造性质，都蕴含邻近纤维具有相同性质。\cmt{V22-DIF015}{此句按原文译出，但“可构造性质”必须指其成立的参数集合可构造。在不可约诺特参数空间中，含泛点的可构造集包含一个非空开集，才得到所说的邻近纤维结论；本版报告给出该拓扑步骤。不能只检查总空间某点的性质，就无条件推出整条纤维、全体中心或任意性质的专门化。}
从记号的角度，应当把$Q$看作（甚至一开始就把它引入为）自然函子态射
\[
\operatorname{Grass}^{m+1}(P)
\longrightarrow
\text{$\operatorname{Grass}^{m}(P)$的子预概形}
\]
的名称。\cmt{V22-DIF016}{法稿先写的“Vari…”被划去，明确手改为“Sous-préschémas”。这里是子预概形函子的语言，未指定一个Hilbert多项式或一个固定Hilbert概形；本版没有替作者增加这些条件。}
